\begin{titlepage}
\centering
\vspace*{3cm}
{\large\textsc{University of the Philippines Cebu}\par}
\vspace{4pt}
{\textsc{Data Structures}\par}
\vspace{2cm}
\rule{0.65\linewidth}{0.4pt}\par\vspace{18pt}
{\Huge\textsc{Second Long Examination}\par}
\vspace{14pt}
{\Large Lecture Notes, Practice Examination\\ and Worked Solutions\par}
\vspace{18pt}
\rule{0.65\linewidth}{0.4pt}\par
\vspace{2cm}
{\large\itshape Array Lists and Vectors \,--\, Singly and Doubly Linked Lists\\[2pt]
Stacks \,--\, Queues \,--\, Deques\par}
\vspace{2cm}
{Compiled from the course lecture notes, slides and sample programs\\
of the Second Long Examination coverage\par}
\vspace{2cm}
{\large D. Juevesano\par}
\vspace{6pt}
{November 2026\par}
\vfill
{\small Every program in this volume was compiled and executed;\\
every printed output is the output of the program shown.\par}
\end{titlepage}

\chapter*{Preface}
\addcontentsline{toc}{chapter}{Preface}

This volume prepares the reader for the Second Long Examination of the course on data
structures. The coverage of the examination is: \emph{lists} (the array list, the
growable array or vector, the singly linked list and the doubly linked list),
\emph{stacks}, \emph{queues} and \emph{deques}. Trees are not part of it. The volume
has three parts.

\textbf{Part I} is a set of consolidated lecture notes. It merges the lecture handouts on
abstract data types, array-based lists and linked lists, the slide deck on the List ADT,
the sample program \cc{arraylist.c} and the linked-list reviewer into one continuous
exposition with definitions, theorems and proofs. Chapters~\ref{ch:stacks},
\ref{ch:queues} and~\ref{ch:deques} treat stacks, queues and deques in full: their
array and linked implementations, the classical applications, and the two techniques
behind most examination problems on them, the \emph{monotonic stack} and the
\emph{monotonic deque}.

\textbf{Part II} is a practice examination written in the manner of the departmental
examination: short conceptual questions, \emph{tracing problems} that carry a twist, and
\emph{problem-solving} items in the style of an online judge (LeetCode, HackerRank),
each with starter code, test cases, the permitted operations, a required time complexity
and, where it applies, a bonus for a faster solution. \textbf{Part III} contains complete
solutions with explanations, diagrams and the usual mistakes.

\para{The rule of abstraction} Every problem states the \emph{permitted operations}. A
solution that reaches into the representation of an ADT --- for instance, indexing the
array inside a stack when only push and pop are permitted --- forfeits the point of the
exercise and, in an examination, usually the credit. Where a problem asks for a new
operation, the solution implements it \emph{on top of} the given interface or, when the
problem says so, directly on the representation (for example, relinking nodes).

\para{Verification} Every listing in this book is included from a source file that was
compiled with \cc{gcc}/\cc{g++} (with all warnings enabled) and run. The outputs printed
in the tracing problems and in Part~III are the real outputs of these programs. Every
solution of Part~III was additionally run against its sample tests and, wherever a simple
reference implementation exists, against that reference on thousands of random inputs;
for the linked structures the tests also count the heap blocks, so that a leaked or
doubly freed node is detected. The two departmental implementation problems
(\emph{eraseItems} and \emph{appendList}) reproduce their sample outputs character for
character.

\para{Conventions} Positions in the lecture-slide array list are 0-based; the sample
program \cc{arraylist.c} uses \emph{1-based} positions in \cc{itemAt} and \cc{setItem}.
Always read the problem statement for the convention in force. Throughout, $n$ denotes
the number of items. The code is plain C (and a little C++): no global variables, no
\cc{static} and no \cc{const}, as in the course.

\para{Sources} The chapter structure follows P.~Morin, \emph{Open Data Structures}
(ArrayStack, ArrayQueue, ArrayDeque, SLList, DLList), which the handouts adapt;
the problem bank is adapted from LeetCode and from the classical exercises of the cited
textbooks. All errors in this volume are the compiler's, or mine.

\tableofcontents
